\appendix

\chapter{StellarOrion Physical Constants and Code Reference}
\label{app:code_constants}

This appendix documents the physical constants and default parameters used in the StellarOrion simulation engine (\texttt{stellarorion\_program\_proc}), providing full traceability between the thesis equations and the Ada/SPARK implementation.

\section{Fundamental Physical Constants}
\label{sec:fundamental_constants}

Table~\ref{tab:constants} lists the physical constants as defined in \texttt{stellarorion\_types.ads}, which are used throughout the aerothermodynamic calculations presented in this thesis.

\begin{table}[H]
\centering
\caption{Fundamental physical constants used in StellarOrion (from \texttt{stellarorion\_types.ads}).}
\label{tab:constants}
\begin{tabular}{llll}
\hline
\textbf{Constant} & \textbf{Symbol} & \textbf{Value} & \textbf{Unit} \\
\hline
Stefan--Boltzmann & $\sigma_B$ & $5.670374419 \times 10^{-8}$ & W\,m$^{-2}$\,K$^{-4}$ \\
Boltzmann & $k_B$ & $1.380649 \times 10^{-23}$ & J\,K$^{-1}$ \\
Avogadro & $N_A$ & $6.02214076 \times 10^{23}$ & mol$^{-1}$ \\
Standard gravity & $g_0$ & $9.80665$ & m\,s$^{-2}$ \\
Molar mass of air & $M_{\text{air}}$ & $28.97 \times 10^{-3}$ & kg\,mol$^{-1}$ \\
Specific gas constant (air) & $R_{\text{air}}$ & $287.058$ & J\,kg$^{-1}$\,K$^{-1}$ \\
Specific heat ratio (air) & $\gamma_{\text{air}}$ & $1.4$ & --- \\
Molecular diameter (air) & $d$ & $3.7 \times 10^{-10}$ & m \\
Prandtl number (air) & $\text{Pr}$ & $0.71$ & --- \\
Sutherland constant & $S$ & $110.4$ & K \\
Sutton--Graves coefficient & $C_{SG}$ & $1.7415 \times 10^{-4}$ & W\,m$^{-2}$\,(kg\,m$^{-3}$)$^{-1/2}$\,(m\,s$^{-1}$)$^{-3}$ \\
Reference viscosity (air) & $\mu_{\text{ref}}$ & $1.716 \times 10^{-5}$ & Pa\,s \\
Sutherland reference temperature & $T_{\text{ref}}$ & $273.15$ & K \\
Specific heat (air, const.\ pressure) & $c_p$ & $1004.0$ & J\,kg$^{-1}$\,K$^{-1}$ \\
\hline
\end{tabular}
\end{table}

\section{Material Survivability Limits}
\label{sec:survivability_limits}

Table~\ref{tab:survivability} lists the material failure thresholds used by the \texttt{Is\_Survivable} predicate in \texttt{stellarorion\_physics.ads} to determine whether a given flight condition exceeds TPS or structural limits.

\begin{table}[H]
\centering
\caption{Material survivability limits enforced by StellarOrion.}
\label{tab:survivability}
\begin{tabular}{lll}
\hline
\textbf{Parameter} & \textbf{Value} & \textbf{Source} \\
\hline
SiC maximum temperature & 2073\,K & Material datasheet / Rapisarda (2023) Sec.\ 4.3 \\
Kapton maximum temperature & 773\,K & DuPont Kapton HN datasheet \\
Maximum g-load & 25\,$g$ & Rapisarda (2023) Sec.\ 5.4 \\
\hline
\end{tabular}
\end{table}

\section{Default Vehicle Geometry (IRVE-3 Baseline)}
\label{sec:vehicle_geometry}

Table~\ref{tab:geometry} presents the default HIAD geometry parameters corresponding to the IRVE-3 flight configuration as implemented in the StellarOrion engine.

\begin{table}[H]
\centering
\caption{Default HIAD geometry parameters (IRVE-3 MDAO baseline from Rapisarda, 2023).}
\label{tab:geometry}
\begin{tabular}{lll}
\hline
\textbf{Parameter} & \textbf{Symbol} & \textbf{Value} \\
\hline
Aeroshell diameter & $D$ & 3.0\,m \\
Cone half-angle & $\theta_c$ & 60\textdegree \\
Nose radius & $R_N$ & 0.55\,m \\
Number of toroids & $N_{\text{tor}}$ & 6 \\
Toroid radius & $r_{\text{tor}}$ & 0.135\,m \\
Outer toroid radius & $r_{\text{out}}$ & 0.1016\,m \\
Vehicle mass & $m$ & 281\,kg \\
Payload height & $h_{\text{pay}}$ & 1.70\,m \\
Slice angle & $\phi$ & 360\textdegree \\
Nose profile & --- & Smooth \\
Skin type & --- & Smooth \\
Scallop points & $N_{\text{sc}}$ & 8 \\
Scallop amplitude & $\delta_{\text{sc}}$ & 0.030\,m \\
\hline
\end{tabular}
\end{table}

\section{Default Flight Conditions}
\label{sec:flight_conditions}

Table~\ref{tab:flight} lists the freestream conditions used as the baseline simulation point, corresponding to an altitude of 52\,km during IRVE-3 reentry.

\begin{table}[H]
\centering
\caption{Default freestream flight conditions at 52\,km altitude.}
\label{tab:flight}
\begin{tabular}{lll}
\hline
\textbf{Parameter} & \textbf{Symbol} & \textbf{Value} \\
\hline
Mach number & $M_\infty$ & 10 \\
Altitude & $h$ & 52\,km \\
Freestream velocity & $V_\infty$ & 2700\,m/s \\
Freestream density & $\rho_\infty$ & $6.9674 \times 10^{-4}$\,kg/m$^3$ \\
Freestream temperature & $T_\infty$ & 270.65\,K \\
\hline
\end{tabular}
\end{table}

\section{Thermal Protection System (TPS) Presets}
\label{sec:tps_presets}

The StellarOrion engine includes several predefined TPS material presets for thermal analysis. Table~\ref{tab:tps} summarizes the available configurations.

\begin{table}[H]
\centering
\caption{TPS material presets available in StellarOrion.}
\label{tab:tps}
\begin{tabular}{llccccc}
\hline
\textbf{Preset} & \textbf{Material} & $\rho$ (kg/m$^3$) & $c_p$ (J/kg\,K) & $k$ (W/m\,K) & $\varepsilon$ & $\delta$ (m) \\
\hline
SiC & Nicalon SiC & 1468 & 1100 & 0.20 & 0.75 & 0.0254 \\
PICA\textsubscript{X} & Phenolic carbon & 320 & 1500 & 0.50 & 0.85 & 0.040 \\
LOFTID & Nextel/Pyrogel & 300 & 1200 & 0.15 & 0.80 & 0.050 \\
Kapton & Polyimide film & 1420 & 1090 & 0.12 & 0.70 & 0.005 \\
Pyrogel & Aerogel insul.\ & 200 & 1000 & 0.02 & 0.85 & 0.025 \\
Multi & Layer composite & 650 & 1050 & 0.10 & 0.80 & 0.040 \\
\hline
\end{tabular}
\end{table}

\section{Computational Chemistry Modes}
\label{sec:chemistry_modes}

Table~\ref{tab:chemistry} documents the gas-phase chemistry models available in the simulation framework.

\begin{table}[H]
\centering
\caption{Chemistry modes implemented in StellarOrion.}
\label{tab:chemistry}
\begin{tabular}{lll}
\hline
\textbf{Mode} & \textbf{Species} & \textbf{Application} \\
\hline
5-Species Earth & N$_2$, O$_2$, NO, N, O & Earth entry (default) \\
11-Species Earth & N$_2$, O$_2$, NO, N, O, N$^+$, O$^+$, NO$^+$, N$_2^+$, O$_2^+$, e$^-$ & High-enthalpy Earth entry with ionization \\
Mars & CO$_2$, CO, N$_2$, O$_2$, NO, C, O, N, CN, CO$^+$, N$_2^+$ & Mars entry \\
\hline
\end{tabular}
\end{table}

\section{Solver Backends}
\label{sec:solver_backends}

Table~\ref{tab:solvers} lists the computational solver backends supported by the framework.

\begin{table}[H]
\centering
\caption{Solver backends in StellarOrion.}
\label{tab:solvers}
\begin{tabular}{lll}
\hline
\textbf{Solver} & \textbf{Type} & \textbf{Description} \\
\hline
SPARTA & DSMC & Primary rarefied gas dynamics solver (Sandia) \\
OpenFOAM & CFD (RANS) & Continuum flow solver for low-altitude conditions \\
PyFluent & CFD (ANSYS) & Commercial CFD interface for parametric studies \\
PyANSYS & CFD (Mechanical) & Structural finite element analysis interface \\
\hline
\end{tabular}
\end{table}

\chapter{StellarOrion Code Architecture Overview}
\label{app:code_architecture}

\section{Module Structure}

The StellarOrion simulation engine is implemented in Ada/SPARK and organized into the following key packages:

\begin{itemize}
    \item \texttt{stellarorion\_types.ads} -- Physical constants, type definitions, vehicle and flight parameter records
    \item \texttt{stellarorion\_physics.ads/adb} -- Aerothermodynamic calculations (mean free path, Knudsen number, Sutton--Graves heat flux, Fay--Riddell correlation, radiative equilibrium, ballistic coefficient, g-load, trajectory profiling)
    \item \texttt{stellarorion\_geometry.ads/adb} -- Parametric CadQuery geometry generation (HIAD, Orion capsule)
    \item \texttt{stellarorion\_simulation.ads/adb} -- SPARTA DSMC simulation driver and Docker integration
    \item \texttt{stellarorion\_optimization.ads/adb} -- Genetic Algorithm optimizer with LHS/CCD design of experiments
    \item \texttt{stellarorion\_pinn.ads/adb} -- PINN refinement via DeepXDE Python sidecar
    \item \texttt{stellarorion\_calibration.ads/adb} -- Automated calibration against IRVE-3 flight data
\end{itemize}

\section{CLI Modes}

The compiled binary supports 21 command-line interface modes including: \texttt{--self-test}, \texttt{--test sample}, \texttt{--compareCalibrate}, \texttt{--optimize}, \texttt{--trajectory}, \texttt{--pinn}, and others. Refer to the StellarOrion README for complete CLI documentation.

\section{GNATprove Verification Status}

All core physics packages (\texttt{stellarorion\_types}, \texttt{stellarorion\_physics}) have been formally verified using GNATprove at proof level 4, ensuring runtime error freedom (overflow, division by zero, index out of bounds) and contract correctness (pre/post-conditions, type invariants).

\chapter{IRVE-3 Validation Data Summary}
\label{app:validation_data}

\section{Peak Aerothermal Loads}

Table~\ref{tab:validation_full} provides the complete validation comparison between the StellarOrion framework predictions and the IRVE-3 post-flight reconstruction data.

\begin{table}[H]
\centering
\caption{Complete IRVE-3 validation comparison (peak values).}
\label{tab:validation_full}
\begin{tabular}{lccc}
\hline
\textbf{Parameter} & \textbf{Flight Data} & \textbf{StellarOrion} & \textbf{Error (\%)} \\
\hline
Peak heat flux ($\dot{q}$) & 14.36\,W/cm$^2$ & 13.83\,W/cm$^2$ & $-3.7$ \\
Total heat load ($Q$) & 195.06\,J/cm$^2$ & 195.17\,J/cm$^2$ & $+0.1$ \\
Peak deceleration & 19.7\,$g$ & 20.2\,$g$ & $+2.5$ \\
Aeroshell diameter & 3.0\,m & 3.0\,m & 0 \\
Toroid count & 6 & 6 & 0 \\
Toroid radius & 0.135\,m & 0.135\,m & 0 \\
\hline
\end{tabular}
\caption*{Source: Flight data from Rapisarda (2023) Table 4.10; StellarOrion values from SPARTA DSMC simulation at grid-factor 0.7.}
\end{table}

\section{Grid Independence Study Results}

Table~\ref{tab:grid_study} summarizes the grid independence study results that established the optimal grid-factor of 0.7.

\begin{table}[H]
\centering
\caption{Grid independence study: stagnation-point heat flux vs.\ grid-factor.}
\label{tab:grid_study}
\begin{tabular}{cc}
\hline
\textbf{Grid-Factor} & \textbf{$\dot{q}_{\text{stag}}$ (W/cm$^2$)} \\
\hline
0.3 & $>$16.0 \\
0.5 & $\sim$14.5 \\
\textbf{0.7} & \textbf{13.83} \\
1.0 & $\sim$14.1 \\
\hline
\end{tabular}
\caption*{The grid-factor of 0.7 provides the minimum error relative to the 14.36\,W/cm$^2$ flight measurement.}
\end{table}

\chapter{Mathematical Derivation Details}
\label{app:derivations}

\section{Sutton--Graves Heat Flux Derivation}

The Sutton--Graves stagnation-point convective heat flux correlation~\cite{sutton1971gravesh} is given by:
\begin{equation}
    \dot{q} = C_{SG} \sqrt{\frac{\rho_\infty}{R_N}} \, V_\infty^3
    \label{eq:sg_derivation}
\end{equation}

\noindent where $C_{SG} = 1.7415 \times 10^{-4}$\,W\,m$^{-2}$\,(kg\,m$^{-3}$)$^{-1/2}$\,(m\,s$^{-1}$)$^{-3}$ is the gas-model-dependent coefficient for a 5-species chemically reacting air model. This correlation captures the $V^3$ velocity dependence arising from the combined effects of shock compression and boundary-layer energy transport.

\subsection{Dimensional Verification}

Substituting the IRVE-3 baseline conditions ($\rho_\infty = 6.9674 \times 10^{-4}$\,kg/m$^3$, $R_N = 0.55$\,m, $V_\infty = 2700$\,m/s):
\begin{align}
    \dot{q} &= 1.7415 \times 10^{-4} \times \sqrt{\frac{6.9674 \times 10^{-4}}{0.55}} \times (2700)^3 \nonumber \\
    &= 1.7415 \times 10^{-4} \times \sqrt{1.2668 \times 10^{-3}} \times 1.9683 \times 10^{10} \nonumber \\
    &= 1.7415 \times 10^{-4} \times 0.03559 \times 1.9683 \times 10^{10} \nonumber \\
    &\approx 12.20 \text{ W/cm}^2
    \label{eq:sg_baseline}
\end{align}

\noindent This single-point value is lower than Rapisarda's trajectory-integrated peak of 15.26\,W/cm$^2$ because the baseline density used here is the steady-state condition at 52\,km, whereas the trajectory-integrated value accounts for the density variation along the full reentry path.

\section{Mean Free Path and Knudsen Number}

The mean free path for a gas of hard-sphere molecules is~\cite{bird1994molecular}:
\begin{equation}
    \lambda = \frac{1}{\sqrt{2} \, \pi \, d^2 \, n}
    \label{eq:mfp_derivation}
\end{equation}

\noindent where $d = 3.7 \times 10^{-10}$\,m is the molecular diameter and $n$ is the number density. The Knudsen number is then:
\begin{equation}
    Kn = \frac{\lambda}{L}
    \label{eq:kn_derivation}
\end{equation}

\noindent where $L$ is the characteristic length (aeroshell diameter, $D = 3.0$\,m).

\section{Fay--Riddell Stagnation-Point Heat Transfer}

The Fay--Riddell equation~\cite{fay1958stagnation} provides the theoretically rigorous stagnation-point heat flux for a reacting, dissociated boundary layer:
\begin{equation}
    \dot{q}_s = 0.763 \, \text{Pr}^{-0.6} \, (\rho_e \mu_e)^{0.4} \, (\rho_w \mu_w)^{0.1} \, \sqrt{\left(\frac{du_e}{dx}\right)_s} \, (h_s - h_w)
    \label{eq:fay_riddell}
\end{equation}

\noindent Unlike the Sutton--Graves correlation, the Fay--Riddell formulation accounts for the actual shock layer structure, chemical non-equilibrium, and catalytic wall effects, making it more physically accurate for HIAD geometries.
